% 来源：解析据 papers/2017-2018-秋冬-期中-甲.tex 撰写（原卷见 sources/2016-2025-期中-甲-历年真题汇编.pdf 第 3 页）

\begin{solution}
记该 $n$ 阶行列式为 $D_n$：主对角线取 $x$，主对角线右上方取 $y$，左下方取 $z$。

\textbf{第一步　把每行化成只有两个非零元。}自下而上作行变换 $R_i\leftarrow R_i-R_{i-1}$（$i=n,n-1,\cdots,2$，
减的是尚未改动的第 $i-1$ 行）：第 $i$ 行只剩第 $i-1$ 列的 $z-x$ 与第 $i$ 列的 $x-y$，行列式不变，得
\[
D_n=\begin{vmatrix}
x&y&y&\cdots&y\\
z-x&x-y&0&\cdots&0\\
0&z-x&x-y&\cdots&0\\
\vdots&\ddots&\ddots&\ddots&\vdots\\
0&\cdots&0&z-x&x-y
\end{vmatrix}=:\det B .
\]

\textbf{第二步　按第一行展开。}$B$ 的第 $i$ 行（$i\ge2$）只在第 $i-1$、$i$ 两列非零。删去第 $1$ 行与第 $k$ 列后，
余子式要非零，必须每行从自己的两个非零元中取一个、且所取列号两两不同。设第 $i$ 行取左元（第 $i-1$ 列，值 $z-x$）的行组成集合 $S$：
若有 $i\in S$、$i\ge3$ 而 $i-1\notin S$，则第 $i-1$ 行取右元仍落在第 $i-1$ 列，与第 $i$ 行所取同列而重复，
故 $S$ 必形如 $\{2,3,\cdots,m\}$；此时所取列号的集合是 $\{1,\cdots,n\}\setminus\{m\}$，而可用列只有 $\{1,\cdots,n\}\setminus\{k\}$，
两者相等给出 $m=k$，即唯一取法 $S=\{2,\cdots,k\}$。所取列号按行号 $2,3,\cdots,n$ 依次为
$1,2,\cdots,k-1,k+1,\cdots,n$，是递增的，故这一项的排列符号为 $+1$。连同展开式自身的 $(-1)^{1+k}$ 得
\[
\det B=x(x-y)^{n-1}+y\sum_{k=2}^{n}(-1)^{k+1}(z-x)^{k-1}(x-y)^{n-k},
\]
其中 $k=1$ 的余子式是上三角阵、对角元全为 $x-y$，值为 $(x-y)^{n-1}$。

\textbf{第三步　求和化简。}记 $u=x-y$、$v=z-x$，上式中的和 $S=\sum_{m=1}^{n-1}(-1)^{m}v^{m}u^{n-1-m}$，把它拆开：
\[
(u+v)S=\sum_{m=1}^{n-1}\Bigl[(-1)^{m}v^{m}u^{n-m}+(-1)^{m}v^{m+1}u^{n-1-m}\Bigr]=-vu^{n-1}+(-1)^{n+1}v^{n},
\]
第 $m$ 项的后一半与第 $m+1$ 项的前一半正负相消，只剩首尾。再由 $u+v=z-y$ 及 $x(z-y)-y(z-x)=z(x-y)$，
\begin{align*}
(z-y)D_n&=x(z-y)u^{n-1}+y\bigl(-vu^{n-1}+(-1)^{n+1}v^{n}\bigr)\\
&=\bigl[x(z-y)-y(z-x)\bigr]u^{n-1}+(-1)^{n+1}yv^{n}
=z(x-y)^{n}+(-1)^{n+1}y(z-x)^{n}.
\end{align*}
这是关于 $x,y,z$ 的多项式恒等式，只有在 $y=z$ 时不能直接除以 $z-y$。

\textbf{情形一：$y\ne z$。}
\begin{align*}
D_n=\frac{z(x-y)^{n}+(-1)^{n+1}y(z-x)^{n}}{z-y}
&=\frac{y(x-z)^{n}-z(x-y)^{n}}{y-z}
\end{align*}
（后一等式用了 $(z-x)^{n}=(-1)^{n}(x-z)^{n}$）。分母中只出现 $z-y$，所以 $x=y$ 或 $x=z$ 都不是退化，
直接代入即可：$x=y$ 时 $D_n=(-1)^{n+1}y(z-y)^{n-1}$，$x=z$ 时 $D_n=z(z-y)^{n-1}$。

\textbf{情形二：$y=z$。}此时所有非对角元都是 $y$，即 $D_n=\bigl|(x-y)E+yJ\bigr|$，$J$ 为全 $1$ 阵。
$J$ 的特征值为 $n$（一重）与 $0$（$n-1$ 重），故
\[
D_n=(x-y)^{n-1}\bigl[(x-y)+ny\bigr]=(x-y)^{n-1}\bigl(x+(n-1)y\bigr).
\]
\ans{$y\ne z$ 时 $D_n=\dfrac{z(x-y)^{n}+(-1)^{n+1}y(z-x)^{n}}{z-y}$（也等于 $\dfrac{y(x-z)^{n}-z(x-y)^{n}}{y-z}$）；
$y=z$ 时 $D_n=(x-y)^{n-1}\bigl(x+(n-1)y\bigr)$。}
\rem{退化只出在 $y=z$（分母为零），此时改用 $M=(x-y)E+yJ$ 的特征值来算；所得结果恰是通项公式在 $z\to y$ 处的极限，
可用来验算。$x=y$、$x=z$ 不属于退化，代入通项即可。}
\end{solution}

\begin{solution}
把第 $i$ 个方程写成 $(k-1)x_i+(x_1+x_2+x_3)=b_i$，其中 $b_1=k-3$、$b_2=b_3=-2$，并记 $s=x_1+x_2+x_3$。
三式相加：左端为 $(k-1)s+3s=(k+2)s$，右端为 $(k-3)-2-2=k-7$，故
\[
(k+2)s=k-7 .
\]
这是全题的关键：$k\ne-2$ 时它定出 $s$，$k=-2$ 时它自相矛盾。

\textbf{$k\ne1,-2$：唯一解。}系数矩阵是 $(k-1)E+J$（$J$ 为全 $1$ 阵），其特征值为 $k-1$（$n-1$ 重）与 $k+2$，
故 $|A|=(k-1)^{2}(k+2)\ne0$，解唯一。由 $(k+2)s=k-7$ 得 $s=\dfrac{k-7}{k+2}$，再由 $(k-1)x_i=b_i-s$ 得
\begin{align*}
x_1&=\frac{1}{k-1}\Bigl[(k-3)-\frac{k-7}{k+2}\Bigr]
=\frac{(k-1)^{2}}{(k-1)(k+2)}=\frac{k-1}{k+2},\\
x_2&=x_3=\frac{1}{k-1}\Bigl[-2-\frac{k-7}{k+2}\Bigr]
=\frac{-3(k-1)}{(k-1)(k+2)}=\frac{-3}{k+2}.
\end{align*}
把它们代回三个原方程即可验算。

\textbf{$k=-2$：无解。}此时 $(k+2)s=k-7$ 变成 $0\cdot s=-9$，矛盾；即 $R(A)=2$ 而增广矩阵 $R(A\mid d)=3$。

\textbf{$k=1$：无穷多解。}三个方程都化为 $x_1+x_2+x_3=-2$，$R(A)=R(A\mid d)=1<3$，通解为
\[
(x_1,x_2,x_3)=(-2-c_1-c_2,\ c_1,\ c_2),\qquad c_1,c_2\ \text{为任意实数}.
\]
\ans{$k\ne1$ 且 $k\ne-2$ 时有唯一解 $x=\Bigl(\dfrac{k-1}{k+2},\ \dfrac{-3}{k+2},\ \dfrac{-3}{k+2}\Bigr)^{T}$；
$k=-2$ 时无解；$k=1$ 时有无穷多解，即 $x_1+x_2+x_3=-2$。}
\end{solution}

\begin{solution}
记 $A=\begin{pmatrix}1&1&1\\0&1&2\\1&0&2\end{pmatrix}$，$B=\begin{pmatrix}1&1&1\\0&0&-6\end{pmatrix}$，则 $X$ 是 $2\times3$ 矩阵。
按第一列展开算 $|A|$：
\[
|A|=1\cdot\begin{vmatrix}1&2\\0&2\end{vmatrix}+1\cdot\begin{vmatrix}1&1\\1&2\end{vmatrix}=2+1=3\ne0,
\]
故 $A$ 可逆，矩阵方程的解唯一，为 $X=BA^{-1}$。

用高斯－若尔当消元求 $A^{-1}$：对 $(A\mid E)$ 作初等行变换化为 $(E\mid A^{-1})$，即
\[
\left(\begin{array}{ccc|ccc}
1&1&1&1&0&0\\
0&1&2&0&1&0\\
1&0&2&0&0&1
\end{array}\right)
\longrightarrow
\left(\begin{array}{ccc|ccc}
1&0&0&\frac23&-\frac23&\frac13\\[1mm]
0&1&0&\frac23&\frac13&-\frac23\\[1mm]
0&0&1&-\frac13&\frac13&\frac13
\end{array}\right),
\]
于是
\begin{align*}
A^{-1}&=\begin{pmatrix}\frac23&-\frac23&\frac13\\[1mm]\frac23&\frac13&-\frac23\\[1mm]-\frac13&\frac13&\frac13\end{pmatrix},\\
X=BA^{-1}&=\begin{pmatrix}1&1&1\\0&0&-6\end{pmatrix}A^{-1}
=\begin{pmatrix}1&0&0\\2&-2&-2\end{pmatrix}.
\end{align*}
验算：$X$ 的第 $1$ 行乘 $A$ 得 $(1,1,1)$，第 $2$ 行乘 $A$ 得 $(0,0,-6)$，正是 $B$ 的两行。
\ans{$X=\begin{pmatrix}1&0&0\\2&-2&-2\end{pmatrix}$。}
\end{solution}

\begin{solution}
思路：取齐次方程组 $Ax=\theta$ 的基础解系充当 $B$ 的列。

由秩－零化度定理（对 $A$ 的 $n$ 列，$\dim\ker A+R(A)=n$），$Ax=\theta$ 的解空间维数为
\[
\dim\ker A=n-R(A)=n-r .
\]
取它的一组基 $\xi_1,\cdots,\xi_{n-r}$，令
\[
B=\bigl(\xi_1,\ \cdots,\ \xi_{n-r},\ \theta,\ \cdots,\ \theta\bigr)_{n\times n},
\]
即前 $n-r$ 列是基础解系、其余 $r$ 列取零列。

因 $A\xi_j=\theta$（$j=1,\cdots,n-r$），$AB$ 的每一列都是零列，故 $AB=O$。又 $\xi_1,\cdots,\xi_{n-r}$ 线性无关，
而 $B$ 的后 $r$ 列全为零，故这 $n-r$ 列构成 $B$ 的列向量组的一个极大无关组，$R(B)=n-r$。
（$r=n$ 时 $n-r=0$，取 $B=O$ 即可。）
\qed
\end{solution}

\begin{solution}
\textbf{(1)} 由 $R(A)=1$ 知 $A\ne O$，故 $A$ 有非零行。取定其中一个非零行，记其为行向量 $a=(a_1,\cdots,a_n)$，则 $a\ne0$。
$R(A)=1$ 意味着行空间是一维的，于是每一行都落在 $a$ 张成的直线上，即存在实数 $b_1,\cdots,b_n$ 使 $A$ 的第 $j$ 行为 $b_ja$。
取定的那一行（设为第 $r$ 行）对应 $b_r=1$，故 $b_1,\cdots,b_n$ 也不全为零。把各行写成一列即
\[
A=\begin{pmatrix}a_1\\ a_2\\ \vdots\\ a_n\end{pmatrix}(b_1,b_2,\cdots,b_n)=(a_1,\cdots,a_n)^{T}(b_1,\cdots,b_n).
\]

\textbf{(2)} 沿用 (1) 的记法 $A=a^{T}b$，其中 $a,b$ 为上述行向量；$ba^{T}=\sum_{i=1}^{n}a_ib_i$ 是一个实数，记为 $k$。
把矩阵乘法的结合律用在中间的一对上：
\[
A^{2}=(a^{T}b)(a^{T}b)=a^{T}(ba^{T})b=k\,a^{T}b=kA .
\]
于是存在实数 $k$ 使 $A^{2}=kA$。（顺便得到 $k=\sum_{i=1}^{n}a_ib_i=\operatorname{tr}A$。）
\qed
\end{solution}

\begin{solution}
把 $G$ 按列分块为 $G=(A\ B\ d\ c)$，用 $C(P)$ 表示矩阵 $P$ 的列空间，则 $\dim C(P)=R(P)$。

$AX=d$ 有解，说明 $d$ 是 $A$ 的列向量组的线性组合，即 $d\in C(A)$。于是 $G$ 的列空间为
\[
C(G)=C(A)+C(B)+\operatorname{span}\{d\}+\operatorname{span}\{c\}=C(A)+C(B)+\operatorname{span}\{c\},
\]
即新增的两列中 $d$ 一列不带来新的维数。由此
\[
R(G)=\dim C(G)\le\dim\bigl(C(A)+C(B)\bigr)+1\le R(A)+R(B)+1,
\]
最后一步用了 $\dim(U+V)\le\dim U+\dim V$；这里 $C(A)+C(B)$ 恰是分块矩阵 $(A\ B)$ 的列空间，其维数为 $R(A\ B)$。

顺带说明：证明只用到 $AX=d$ 有解。$BX=c$ 无解等价于 $c\notin C(B)$，即 $R(B\ c)=R(B)+1$，
这说明多出的那一维确由 $c$ 贡献、上界中的 $1$ 一般不能更小：例如取 $A=O_{m\times n}$、$d=\theta$、$B=O_{m\times s}$、$c\ne\theta$
（此时 $AX=d$ 有解、$BX=c$ 无解），有 $R(G)=1=R(A)+R(B)+1$。
\qed
\end{solution}

\begin{solution}
记 $M=\begin{pmatrix}A&D\\C&B\end{pmatrix}$。先对可逆的 $A$ 证，再用摄动法去掉可逆性。

\textbf{第一步：$A$ 可逆时用分块消元。}由 $AC=CA$ 可作如下分块行变换（左乘一个分块初等矩阵）：
\[
\begin{pmatrix}E&O\\-C&A\end{pmatrix}\begin{pmatrix}A&D\\C&B\end{pmatrix}
=\begin{pmatrix}A&D\\-CA+AC&-CD+AB\end{pmatrix}
=\begin{pmatrix}A&D\\O&AB-CD\end{pmatrix},
\]
左下角因 $AC=CA$ 而为零。左边第一个因子是分块三角阵，其行列式为 $|E|\cdot|A|=|A|$，对等式两边取行列式得
\[
|A|\cdot|M|=|A|\cdot|AB-CD| .
\]
$A$ 可逆时 $|A|\ne0$，约去 $|A|$ 即得 $|M|=|AB-CD|$。

\textbf{第二步：$A$ 不可逆时用摄动法。}对参数 $t$ 令 $A(t)=A+tE$。由 $AC=CA$ 得
\[
A(t)C=AC+tC=CA+tC=CA(t),
\]
即 $A(t)$ 仍与 $C$ 可交换，故可对 $A(t)$ 套用第一步的结论。而 $|A(t)|=|A+tE|$ 是 $t$ 的 $n$ 次首一多项式，
至多有 $n$ 个根，所以除有限多个 $t$ 外 $A(t)$ 都可逆。于是
\[
f(t):=\begin{vmatrix}A(t)&D\\C&B\end{vmatrix}=\bigl|A(t)B-CD\bigr|=:g(t)
\]
对无穷多个 $t$ 成立。$f(t)$ 与 $g(t)$ 都是 $t$ 的多项式（行列式按定义展开即若干元素乘积之和），
两个多项式在无穷多个点取值相同，故是同一个多项式；令 $t=0$ 便得
\[
\begin{vmatrix}A&D\\C&B\end{vmatrix}=|AB-CD| .
\]
\qed
\rem{“可逆时消元、不可逆时摄动”是分块行列式恒等式的通用套路：恒等式两边都是矩阵元素的多项式，
只要在可逆（稠密）的情形验证，就能自动延拓到奇异情形，不必对奇异矩阵单独另找消元方案。}
\end{solution}
